% GENERATED FILE. Edit canonical lessons, then run npm run generate:books.
% canonical-source-hash: fnv1a64:47e0c7d064cce052
% canonical-lessons: MR-C498-balpan, MR-C498-tarunya, MR-C498-vardhakya, MR-C498-vay, MR-C498-kisor

\chapter{Growing Up and Growing Old}
\label{ch:mr-growing-up-and-growing-old}
\hlchaptermodality{498}

\begin{chapteropening}
\textbf{By the end of this chapter:} \emph{I can talk about childhood, youth, old age, age and the teenage years.}

Complete the last lesson of chapter 498: I can talk about childhood, youth, old age, age and the teenage years.
\end{chapteropening}

\section[\texorpdfstring{\mr{बालपण} (bālpaṇ)}{बालपण (bālpaṇ)}]{\mr{बालपण} (bālpaṇ) — childhood}

\label{lesson:MR-C498-balpan}



\begin{quote}
Before the new one: say the Marathi for of a paternal uncle, then the Marathi for step-.
\end{quote}

\subsection*{You'll want to know: \mr{बालपण}}
\textbf{\mr{बालपण}} — \emph{bālpaṇ} — \textquotedblleft{}childhood\textquotedblright{}.

\begin{quote}
\mr{माझं} \mr{बालपण} \mr{गावात} \mr{गेलं}.
\end{quote}

\begin{quote}
\emph{mājhaṁ bālpaṇ gāvāt gelaṁ.}
\end{quote}

\begin{quote}
I spent my childhood in the village.
\end{quote}

Neuter.

\mr{बाल}, a child, + -\mr{पण}, -hood.

\begin{cousinweb}
\textbf{\mr{बालपण}} (\emph{bālpaṇ}) is a Marathi derivation: \mr{बाल}, \textquotedblleft{}child\textquotedblright{}, a tatsama of Sanskrit \emph{bāla} (CDIAL 9216), with \emph{-paṇ}, \textquotedblleft{}-hood\textquotedblright{}.
\end{cousinweb}

\begin{grammarlens}[title={What you've built}]
One more word for describing what you lived and felt.
\end{grammarlens}

\subsection*{Your turn}
\begin{itemize}
\raggedright
  \item \emph{Say it:} \emph{bālpaṇ}
  \item \emph{Say it:} \emph{bālpaṇ}, once more
  \item \emph{Say it:} \emph{mājhaṁ bālpaṇ gāvāt gelaṁ.} — I spent my childhood in the village.
  \item \emph{Recall:} say the Marathi for of a paternal uncle, then the Marathi for step-, then say \emph{bālpaṇ} again
\end{itemize}

\begin{tcolorbox}[breakable,colback=teal!4,colframe=teal!35!black,title={Before you move on}]
What does \mr{बालपण} mean? (\textquotedblleft{}Childhood\textquotedblright{}.) Say it once more.
\end{tcolorbox}
\section[\texorpdfstring{\mr{तारुण्य} (tāruṇya)}{तारुण्य (tāruṇya)}]{\mr{तारुण्य} (tāruṇya) — youth}

\label{lesson:MR-C498-tarunya}



\begin{quote}
Before the new one: say the Marathi for step-, then the Marathi for childhood.
\end{quote}

\subsection*{You'll want to know: \mr{तारुण्य}}
\textbf{\mr{तारुण्य}} — \emph{tāruṇya} — \textquotedblleft{}youth\textquotedblright{}.

\begin{quote}
\mr{तारुण्यात} \mr{त्याने} \mr{खूप} \mr{प्रवास} \mr{केला}.
\end{quote}

\begin{quote}
\emph{tāruṇyāt tyāne khūp pravās kelā.}
\end{quote}

\begin{quote}
He travelled a lot in his youth.
\end{quote}

Neuter.

From \mr{तरुण}, young.

\begin{cousinweb}
\textbf{\mr{तारुण्य}} (\emph{tāruṇya}) is a tatsama, taken from Sanskrit \emph{tāruṇya}, \textquotedblleft{}youth, youthfulness\textquotedblright{} (Monier-Williams).
\end{cousinweb}

\begin{grammarlens}[title={What you've built}]
One more word for describing what you lived and felt.
\end{grammarlens}

\subsection*{Your turn}
\begin{itemize}
\raggedright
  \item \emph{Say it:} \emph{tāruṇya}
  \item \emph{Say it:} \emph{tāruṇya}, once more
  \item \emph{Say it:} \emph{tāruṇyāt tyāne khūp pravās kelā.} — He travelled a lot in his youth.
  \item \emph{Recall:} say the Marathi for step-, then the Marathi for childhood, then say \emph{tāruṇya} again
\end{itemize}

\begin{tcolorbox}[breakable,colback=teal!4,colframe=teal!35!black,title={Before you move on}]
What does \mr{तारुण्य} mean? (\textquotedblleft{}Youth\textquotedblright{}.) Say it once more.
\end{tcolorbox}
\section[\texorpdfstring{\mr{वार्धक्य} (vārdhakya)}{वार्धक्य (vārdhakya)}]{\mr{वार्धक्य} (vārdhakya) — old age}

\label{lesson:MR-C498-vardhakya}



\begin{quote}
Before the new one: say the Marathi for childhood, then the Marathi for youth.
\end{quote}

\subsection*{You'll want to know: \mr{वार्धक्य}}
\textbf{\mr{वार्धक्य}} — \emph{vārdhakya} — \textquotedblleft{}old age\textquotedblright{}.

\begin{quote}
\mr{वार्धक्यात} \mr{आजोबा} \mr{गावी} \mr{शांतपणे} \mr{राहतात}.
\end{quote}

\begin{quote}
\emph{vārdhakyāt ājobā gāvī śāntpaṇe rāhtāt.}
\end{quote}

\begin{quote}
In his old age Grandfather lives quietly in the village.
\end{quote}

Neuter.

Formal, beside \mr{म्हातारपण}.

\begin{cousinweb}
\textbf{\mr{वार्धक्य}} (\emph{vārdhakya}) is a tatsama, taken from Sanskrit \emph{vārddhakya}, \textquotedblleft{}old age\textquotedblright{} (Monier-Williams).
\end{cousinweb}

\begin{grammarlens}[title={What you've built}]
One more word for describing what you lived and felt.
\end{grammarlens}

\subsection*{Your turn}
\begin{itemize}
\raggedright
  \item \emph{Say it:} \emph{vārdhakya}
  \item \emph{Say it:} \emph{vārdhakya}, once more
  \item \emph{Say it:} \emph{vārdhakyāt ājobā gāvī śāntpaṇe rāhtāt.} — In his old age Grandfather lives quietly in the village.
  \item \emph{Recall:} say the Marathi for childhood, then the Marathi for youth, then say \emph{vārdhakya} again
\end{itemize}

\begin{tcolorbox}[breakable,colback=teal!4,colframe=teal!35!black,title={Before you move on}]
What does \mr{वार्धक्य} mean? (\textquotedblleft{}Old age\textquotedblright{}.) Say it once more.
\end{tcolorbox}
\section[\texorpdfstring{\mr{वय} (vay)}{वय (vay)}]{\mr{वय} (vay) — age}

\label{lesson:MR-C498-vay}



\begin{quote}
Before the new one: say the Marathi for youth, then the Marathi for old age.
\end{quote}

\subsection*{You'll want to know: \mr{वय}}
\textbf{\mr{वय}} — \emph{vay} — \textquotedblleft{}age\textquotedblright{}.

\begin{quote}
\mr{तुझं} \mr{वय} \mr{किती} \mr{आहे}?
\end{quote}

\begin{quote}
\emph{tujhaṁ vay kitī āhe?}
\end{quote}

\begin{quote}
How old are you?
\end{quote}

Neuter.

Neuter: \mr{माझं} \mr{वय} \mr{दहा} \mr{वर्षं} \mr{आहे}. \mr{वयस्कर}, elderly.

\begin{cousinweb}
\textbf{\mr{वय}} (\emph{vay}) is a tatsama, taken from Sanskrit \emph{vayas}, \textquotedblleft{}energy, prime of life, age\textquotedblright{} (CDIAL 11305), without its final \emph{-s}.
\end{cousinweb}

\begin{grammarlens}[title={What you've built}]
One more word for describing what you lived and felt.
\end{grammarlens}

\subsection*{Your turn}
\begin{itemize}
\raggedright
  \item \emph{Say it:} \emph{vay}
  \item \emph{Say it:} \emph{vay}, once more
  \item \emph{Say it:} \emph{tujhaṁ vay kitī āhe?} — How old are you?
  \item \emph{Recall:} say the Marathi for youth, then the Marathi for old age, then say \emph{vay} again
\end{itemize}

\begin{tcolorbox}[breakable,colback=teal!4,colframe=teal!35!black,title={Before you move on}]
What does \mr{वय} mean? (\textquotedblleft{}Age\textquotedblright{}.) Say it once more.
\end{tcolorbox}
\section[\texorpdfstring{\mr{किशोर} (kiśor)}{किशोर (kiśor)}]{\mr{किशोर} (kiśor) — a teenager}

\label{lesson:MR-C498-kisor}



\begin{quote}
Before the new one: say the Marathi for old age, then the Marathi for age.
\end{quote}

\subsection*{You'll want to know: \mr{किशोर}}
\textbf{\mr{किशोर}} — \emph{kiśor} — \textquotedblleft{}a teenager\textquotedblright{}.

\begin{quote}
\mr{या} \mr{मासिकात} \mr{किशोरांसाठी} \mr{गोष्टी} \mr{असतात}.
\end{quote}

\begin{quote}
\emph{yā māsikāt kiśorāṁsāṭhī goṣṭī astāt.}
\end{quote}

\begin{quote}
This magazine has stories for teenagers.
\end{quote}

Masculine.

\mr{किशोरी}, a teenage girl; \mr{किशोरवय}, the teenage years.

\begin{cousinweb}
\textbf{\mr{किशोर}} (\emph{kiśor}) is a tatsama, taken from Sanskrit \emph{kiśora}, \textquotedblleft{}a foal; a youth\textquotedblright{} (CDIAL 3190).
\end{cousinweb}

\begin{grammarlens}[title={What you've built}]
One more word for describing what you lived and felt.
\end{grammarlens}

\subsection*{Your turn}
\begin{itemize}
\raggedright
  \item \emph{Say it:} \emph{kiśor}
  \item \emph{Say it:} \emph{kiśor}, once more
  \item \emph{Say it:} \emph{yā māsikāt kiśorāṁsāṭhī goṣṭī astāt.} — This magazine has stories for teenagers.
  \item \emph{Recall:} say the Marathi for old age, then the Marathi for age, then say \emph{kiśor} again
\end{itemize}

\begin{tcolorbox}[breakable,colback=teal!4,colframe=teal!35!black,title={Before you move on}]
What does \mr{किशोर} mean? (\textquotedblleft{}A teenager\textquotedblright{}.) Say it once more.
\end{tcolorbox}
